\sectionbrief
{a scalar equation in one space variable, $u_t=\sum_nT_n(u)$, each $T_n$ linear or polynomial in the jet $(u,u_x,u_{xx},\dots)$; $\Omega$ bounded and simply connected; the exact solution is known.}
{to test which of the notions of Sections~\ref{sec:ident}--\ref{sec:inside} survive when the terms are nonlinear.}
{extension of the chain. \emph{Result:} weaker, relative statements (exact expansion of $\Eint$, sharp asymptotics across a shock for Burgers, identification relative to a dictionary). Benchmarks: Burgers, KdV--Burgers, Fisher--KPP; KdV and NLS only as controls.}
We consider $u_t=\sum_nT_n(u)$ with each $T_n$ linear or polynomial in the jet $(u,u_x,u_{xx},\dots)$. Benchmarks: Burgers $u_t=-uu_x+\nu u_{xx}$, KdV--Burgers, Fisher--KPP; integrable cases (KdV, NLS) serve as controls.

\subsection{Interaction term for flows of single terms}
With $S_n$ the (nonlinear) flow of $u_t=T_n(u)$ alone, \eqref{xeq:eint-def} still defines $\Eint^{\mathrm{add}}$. Its Taylor expansion in $t$:
\begin{theorem}[{(Jet expansion) \led{S4T-T1} \proved[]}]\label{thm:jet}
$\displaystyle \Eint^{\mathrm{add}}=\frac{t^2}{2}\sum_{n\ne m}T_n'[\,T_m(u_0)\,]+O(t^3)$, where $T_n'[v]$ is the Fr\'echet derivative of $T_n$ at $u_0$ in the direction $v$.
\end{theorem}
\begin{example}[Burgers]\label{ex:burgers-jet}
$T_1=-uu_x$, $T_2=\nu u_{xx}$. Then $T_1'[v]=-(vu_x+uv_x)$ and $T_2'[v]=\nu v_{xx}$, so
\[
\begin{aligned}
T_1'[T_2u_0]&=-\nu\,(u_0''u_0'+u_0u_0'''),\\
T_2'[T_1u_0]&=-\nu\,(3u_0'u_0''+u_0u_0'''),
\end{aligned}
\]
and half the sum times $t^2$ gives
\[
\begin{aligned}
\Eint&=\tfrac12\big(T_1'[T_2u_0]+T_2'[T_1u_0]\big)t^2+O(t^3)\\
&=-\nu\big(u_0u_0'''+2u_0'u_0''\big)t^2+O(t^3).
\end{aligned}
\]
\end{example}

\subsection{Before, at, and after the shock}
The interaction term of Burgers changes qualitatively with time, and each regime calls for different tools:
\begin{description}[nosep,leftmargin=3.2cm,labelwidth=3.0cm]
\item[Before shock formation] smoothness permits perturbative expansions and derivative-based estimates.
\item[At breaking] the inviscid gradient blows up, invalidating estimates that require smoothness.
\item[After shock formation] pointwise behavior can be discontinuous, so $L^1$ estimates become more informative than pointwise bounds.
\end{description}
Figure~\ref{fig:burgers} recalls why a shock appears at all in the flow of the single term $-uu_x$: along each characteristic $u$ is constant and the characteristic moves with speed $u$, so faster parts overtake slower ones. For $u_0=\sin x$ the characteristics are
\[
x=\xi+t\sin\xi
\]
(the parameter $s$ in the figure plays the role of the time $t$); since
\[
\frac{\partial x}{\partial\xi}=1+t\cos\xi
\]
vanishes first at $\xi=\pi$, the breaking time is $t^*=1$.
\begin{figure}[!htb]
\centering
\begin{tikzpicture}
\begin{axis}[width=9.5cm,height=5.0cm,xlabel={$x$},ylabel={$u$},xmin=0,xmax=6.283,ymin=-1.9,ymax=1.9,
 legend style={font=\scriptsize,draw=none,fill=white,fill opacity=0.85,text opacity=1,at={(0.01,0.02)},anchor=south west},
 title={\small Inviscid Burgers, $u_0=\sin x$: $x=\xi+s\sin\xi,\ u=\sin\xi$},no markers,domain=0:6.283185,samples=200,smooth]
\addplot[cG,thick,dashed,variable=\z] ({\z},{sin(deg(\z))});
\addplot[cA,thick,variable=\z] ({\z+0.5*sin(deg(\z))},{sin(deg(\z))});
\addplot[cC,thick,variable=\z] ({\z+1.0*sin(deg(\z))},{sin(deg(\z))});
\addplot[cB,thick,variable=\z] ({\z+1.6*sin(deg(\z))},{sin(deg(\z))});
\draw[cB,dashed] (axis cs:3.14159,-1) -- (axis cs:3.14159,1);
\legend{$s=0$,$s=0.5$,$s=1$ (breaking),$s=1.6$ (multivalued)}
\end{axis}
\end{tikzpicture}
\caption{Characteristic solution of the inviscid flow of $-uu_x$. At $s=1$ the profile has a vertical tangent: the gradient blows up. Beyond that the curve folds; the physical (entropy) solution replaces the fold by a jump at $x=\pi$ (equal-area rule), shown dashed. This breaking is what makes $\Eint$ discontinuous after the shock.}\label{fig:burgers}
\end{figure}
\paragraph{Before the shock ($t<t^*$): smooth regime.} Smoothness permits perturbative expansions and derivative-based estimates.
\begin{itemize}[nosep]
\item $\Eint$ has a closed form at first order in $\nu$, with remainder $O(\nu^2\varepsilon^2t^3)$ \led{S4T-T2} \proved; it blows up at the breaking time $t^*$ \led{S4T-P3}. (This is the same family of expansions as Theorem~\ref{thm:jet} and Example~\ref{ex:burgers-jet}.)
\item The Cole--Hopf substitution $u=-2\nu\varphi_x/\varphi$ turns Burgers into the heat equation $\varphi_t=\nu\varphi_{xx}$; the inverse scattering transform linearizes KdV and NLS. These are used as the ``exact linearization'' of the nonlinear solution~\cite{ch}.
\end{itemize}
\paragraph{At breaking ($t=t^*$): estimates that need smoothness fail.} The inviscid gradient $u_x$ blows up (vertical tangent, $s=1$ in Figure~\ref{fig:burgers}), so every statement that needs bounds on $u_x,u_{xx},\dots$ stops being usable at $t^*$: the closed form of $\Eint$ blows up, and estimates that require smoothness, such as those behind the jet expansion of Theorem~\ref{thm:jet}, are no longer valid.
\paragraph{After the shock ($t>t^*$): discontinuous regime.}
The flow of $-uu_x$ alone forms a shock, so $\Eint$ jumps by amounts independent of $\nu$; this is forced by the definition. A function with a jump is not in $H^1$ in one dimension, so the Sobolev norms of Sections~\ref{sec:groups} and~\ref{sec:inside} are no longer the natural ones: pointwise values are unreliable, while the \emph{measure} of the discrepancy is well behaved. The meaningful quantities are therefore $L^1$ sizes, collected below:
\begin{align}
\|\Eint\|_{L^1}&=\Theta(\nu),\\
\sup_{t\le T}\|\Eint\|_{L^1}&=\nu|\log\nu|+O(\nu),\\
\|S-S_1S_2\|_{L^1}&=\Theta(\sqrt\nu).
\end{align}
The first line holds at fixed $t$, each shock contributing $4\nu\log\!\big(2\cosh(\rho/2)\big)$; the second has exact leading coefficient $1$. Sources: the first line is \led{OPA-T1} \modulo{Kru\v zkov uniqueness~\cite{kr}}; the second is \led{OPF5A-T1/T2}; the third is \led{OPF5D-T6} \proved.
The second line holds for data with a single birth of shock such as $\sin x$, and is conditional on a localization lemma in general \condl{}.

\begin{table}[!htb]
\centering\small
\begin{tabular}{@{}L{3.0cm}L{4.9cm}L{6.2cm}@{}}
\toprule
Regime & What happens & Tools and results\\\midrule
Before the shock, $t<t^*$ & the profile is smooth, characteristics do not cross & jet expansion (Thm.~\ref{thm:jet}), first-order closed form of $\Eint$, Cole--Hopf\\
At breaking, $t=t^*$ & the inviscid gradient blows up (vertical tangent) & smoothness-based estimates fail; $\Eint$ blows up\\
After the shock, $t>t^*$ & the single-term flow is discontinuous (equal-area rule) & $L^1$ estimates: $\Theta(\nu)$, $\nu|\log\nu|$, $\Theta(\sqrt\nu)$\\
\bottomrule
\end{tabular}
\caption{The three regimes of the Burgers interaction term and the tools that apply in each.}\label{tab:regimes}
\end{table}

\subsection{Identification relative to a dictionary}
Nonlinear terms cannot be handled by the ideal $\Ann(u)$ of Section~\ref{sec:theory}; instead one fixes a finite candidate list. Choose a dictionary $\mathcal D$ of monomials of degree $\le d$ in the jet, of order $\le k$, and ask which polynomial relations $F(u,u_x,\dots)=0$ the data satisfy. Notation: $V$ is the span of $\mathcal D$, $I(u)$ the ideal of all differential relations satisfied by $u$, $[F]$ the differential ideal generated by $F$, and $\mathcal J_u$ the jet-level obstruction (a different object from the cofactor $J_u$ of Section~\ref{sec:theory}, hence the calligraphic letter)\edmark{E13}.
\begin{theorem}[{(Dictionary criterion) \led{S4ID-T1} \proved[]}]\label{thm:dict}
For smooth $u$ on any $\Omega$ there is an exact sequence $0\to[F]\cap V\to I(u)\cap V\to\mathcal J_u\cap\psi(V)\to0$. In particular
\[
I(u)\cap V=[F]\cap V\iff\mathcal J_u\cap\psi(V)=0 .
\]
\end{theorem}
The correct statement is therefore ``the relations are exactly $F$ \emph{modulo its differential consequences}''; the literal ``$I(u)\cap V=\langle F\rangle$'' is false in complete dictionaries (they contain $\partial_xF$, etc.) \led{S4ID-P1}. Concrete criteria: for Burgers with $N$ shocks, $u$ identifies the equation in the evolutionary dictionary iff $d\le N$ \led{S4ID-T2, OP-T7}; for $N$-soliton KdV the identifying polynomial has total degree $(N+2)2^{N-1}$ (generic $N=2$: $8$) \led{OP-K15} \modulo{Novikov 1974~\cite{no}}. Finally, ``closed form $\Rightarrow$ does not identify'' is itself a theorem \led{S4ID-T4}.
\edmark{E2}

\subsection{Linear versus nonlinear: the same idea, two structures}
Sections~\ref{sec:theory}--\ref{sec:inside} and the present one answer the same question: \emph{what does a single solution $u$ remember about its equation?} In both settings we collect all relations satisfied by $u$ and ask when these relations are generated by the equation we were given. What changes is the algebraic structure in which the relations live.
\begin{table}[!htb]
\centering\small
\begin{tabular}{@{}L{2.8cm}L{6.2cm}L{6.2cm}@{}}
\toprule
 & Linear approach (Sections~\ref{sec:theory}--\ref{sec:inside}) & Nonlinear approach (this section)\\\midrule
Space of candidate equations & the ring $\ring=\C[\tau,\xi]$ of constant-coefficient operators & a finite \emph{dictionary} $\mathcal D$ of jet monomials of degree $\le d$ and order $\le k$; $V=\spn\mathcal D$\\
Relations satisfied by $u$ & $\Ann(u)$, an ideal of $\ring$ & $I(u)$, the ideal of all differential relations\\
What the equation generates & the principal ideal $(P)$ & the differential ideal $[F]$\\
Criterion for identification & $\Ann(u)=(P)$ (Theorem~\ref{thm:ident}) & $I(u)\cap V=[F]\cap V$, i.e.\ $F$ \emph{modulo its differential consequences} (Theorem~\ref{thm:dict})\\
What escapes the equation & the cofactor $J_u$: the finitely many exponentials in $u$ ($\sigma_p(u)=\Var(J_u)$) & the jet-level obstruction $\mathcal J_u\cap\psi(V)$ (exact sequence of Theorem~\ref{thm:dict})\\
Solutions that do not identify & $u\in\EP$ & closed-form solutions (\led{S4ID-T4})\\
Concrete criteria & finite linear algebra (Lemma~\ref{lem:ep}) & Burgers with $N$ shocks identifies iff $d\le N$; $N$-soliton KdV: degree $(N+2)2^{N-1}$\\
\bottomrule
\end{tabular}
\caption{The linear and nonlinear identification problems side by side.}\label{tab:lin-nonlin}
\end{table}
\begin{aside}{Where the parallel breaks}
\textbf{Linear case.} The solutions form a vector space and the annihilators an ideal: one can add, multiply by operators and factor, which is why the Nullstellensatz applies and the test is a finite-dimensional computation.

\textbf{Nonlinear case.} There is no linear structure on the solutions: one must fix a finite space $V$ of candidates in advance, and the identification is \emph{relative to $V$}, never absolute. The literal statement
\[
I(u)\cap V=\langle F\rangle
\]
is false in complete dictionaries, which contain $\partial_xF$, so the correct statement is ``modulo differential consequences''.
\end{aside}

\begin{keypoint}
Both approaches ask what one solution determines about its equation. The linear one answers with the annihilator ideal $\Ann(u)$ and its generator $G_u$; the nonlinear one answers with a dictionary and the differential ideal of relations, and only up to differential consequences. In the algorithm of the companion paper Steps~1--6 use the structure of $\Ann(u)$, while Step~7 switches to the dictionary and returns a relative identification.
\end{keypoint}

%======================================================================
